\documentclass[11pt,a4paper,twocolumn]{article}
\usepackage[utf8]{inputenc}
\usepackage[T1]{fontenc}
\usepackage{amsmath, amssymb, graphicx, booktabs, tikz}
\usepackage{microtype}
\usepackage[hyphens]{url}
\usetikzlibrary{arrows.meta, positioning, shapes.geometric}
\raggedbottom

\title{The Alpha Difficulty Scoring Model}
\author{}
\date{\today}

\begin{document}

\maketitle

\begin{abstract}
Alpha is a first, intentionally simple model for scoring the navigational
difficulty of a route through an indoor orienteering venue. It is additive
over two structural quantities --- stair crossings and branching points ---
and its scale is calibrated so that scores align with the difficulty
categories published by Stockholm Indoor
Cup~\cite{sic-difficulty}. This note defines the model in isolation, as a
baseline against which future scoring models can be compared; it is a
condensed excerpt of the difficulty-scoring model introduced in the
companion article on graph-based indoor orienteering
analysis~\cite{foundation}.
\end{abstract}

\section{Model definition}
\label{sec:model}

Alpha assumes a route has already been identified (as a sequence of graph
edges) and scores it from two counts:

\[
D \;=\; k \bigl(5\,N_f + 3\,N_b\bigr), \qquad k = \frac{100}{81}
\]

where

\begin{itemize}
    \item $N_f$ is the number of stair (floor-change) traversals along the
          route, upward and downward combined --- an edge counts as a stair
          when its endpoints resolve to different floor levels;
    \item $N_b$ is the number of \textbf{branching points}: route nodes at
          which more than one next node is reachable, excluding the node
          just come from and respecting one-way constraints;
    \item $k = 100/81$ is a scaling factor that normalizes a reference leg
          (Stockholm Indoor Cup 2023, stage 1, course 1, leg 3$\to$4; raw
          score 81) to exactly 100.
\end{itemize}

Leg difficulty is the score of the leg's shortest route; course difficulty
defaults to the \emph{maximum} leg difficulty. Scores map onto the five-color
scale at fixed thresholds (Table~\ref{tab:thresholds}). Alpha is implemented
as part of the open-source analysis toolchain~\cite{repo}, which computes it
directly from a venue graph.

\begin{table}[h]
\centering
\begin{tabular}{@{}ll@{}}
\toprule
\textbf{Category} & \textbf{Score range} \\ \midrule
White  & $0$--$20$  \\
Yellow & $20$--$40$ \\
Orange & $40$--$60$ \\
Purple & $60$--$80$ \\
Black  & $80+$      \\ \bottomrule
\end{tabular}
\caption{Difficulty scores mapped to difficulty categories.}
\label{tab:thresholds}
\end{table}

\section{Worked example}
\label{sec:example}

Figure~\ref{fig:example} shows a two-floor fragment with a single route from
start $S$ to finish $F$. The route crosses two labeled stairs (A and B) and
passes through one branching point, where a dead-end corridor is also
reachable.

\begin{figure}[h]
\centering
\begin{tikzpicture}[every node/.style={draw, circle, inner sep=1.8pt},
    label/.style={draw=none, inner sep=1pt, font=\footnotesize}]
    % Floor 1 (bottom)
    \node (S)  at (0,0)      {};
    \node (B1) at (1.4,0)    {};
    \node (DE) at (2.4,-1.0) {};
    \node (F)  at (3.6,0)    {};
    \draw (S) -- (B1);
    \draw (B1) -- (DE);

    % Floor 2 (top)
    \node (M) at (1.4,1.6) {};
    \node (N) at (3.6,1.6) {};
    \draw (M) -- (N);

    % Stairs (dashed, per graph notation)
    \draw [dashed] (B1) -- (M) node [label, midway, left] {A};
    \draw [dashed] (N) -- (F) node [label, midway, right] {B};

    % Divider between floor 1 and floor 2
    \draw [gray, dotted] (-1.1,0.8) -- (4.0,0.8);

    % Node and floor labels (floor labels share a common x so they align)
    \node [label, above=2pt] at (S)  {$S$};
    \node [label, below=2pt] at (F)  {$F$};
    \node [label, below=2pt] at (DE) {dead end};
    \node [label] at (-0.9,0)   {floor 1};
    \node [label] at (-0.9,1.6) {floor 2};
\end{tikzpicture}
\caption{Worked example. The route from $S$ to $F$ crosses stairs A and B
($N_f = 2$) via one branching point --- the node where the dead-end corridor
also departs ($N_b = 1$). This gives
$D = k(5 \cdot 2 + 3 \cdot 1) = \frac{100}{81}\cdot 13 \approx 16.0$, in the
White category.}
\label{fig:example}
\end{figure}

\section{Refinements beyond raw branching}
\label{sec:refinements}

Raw branching counts every fork, including trivially dismissible ones (such
as the dead end in Figure~\ref{fig:example}, which is short and unlikely to
mislead). Two refinements are computed alongside the score and available to
future models:

\begin{itemize}
    \item \textbf{Dead-end filtering:} edges on non-branching chains that
          terminate at a leaf are excluded --- a passage with only one way in
          and out creates no real decision.
    \item \textbf{Decision points:} a stricter variant that counts a node
          only when an off-route branch is \emph{significant}, i.e.\ runs at
          least 50 m before dead-ending or rejoining the route. A 10 m alcove
          is visibly wrong and creates no decision; a 50 m parallel corridor
          is a genuine route choice.
\end{itemize}

\section{Known limitations}
\label{sec:limitations}

Alpha is a baseline, and its limitations chart the research agenda:
branching is counted, not weighted by how hard the alternatives are to
reject; the additive form cannot express the combinatorial explosion of
several \emph{sequential} decisions; and the score depends on how densely
the navigation graph is drawn, so graph-drawing conventions must be
standardized before scores are comparable across venues. These are taken up
as research directions in the companion article~\cite{foundation}.

\begin{thebibliography}{9}

\bibitem{foundation}
\emph{A graph-based framework for quantitative analysis of navigational
complexity in indoor orienteering courses} (companion article).

\bibitem{sic-difficulty}
Stockholm Indoor Cup. \emph{Sv\r{a}righetsgrader --- Difficulty levels}.

\bibitem{repo}
\emph{Indoor Orienteering Route Analysis} (source code, algorithm reference,
and datasets). \newblock
\url{github:objorke/indoor-o-route-analysis}.

\end{thebibliography}

\end{document}
